\chapter{Sécurité, tests et qualité}
\label{ch:securite}

\section{Modèle de menace (périmètre local)}

Le déploiement documenté cible un \textbf{environnement local de démonstration}. Les menaces principales sont : accès non autorisé à l'API, fuite de secrets dans Git, injection SQL (mitigée par ORM), et configuration CORS trop permissive.

\section{Authentification et autorisation}

\begin{itemize}
  \item Mots de passe hashés bcrypt — pas de stockage clair ;
  \item JWT signé HS256 — secret \texttt{SECRET\_JWT} injecté par variable d'environnement ou Secret K8s ;
  \item Durée de vie configurable (\texttt{duree\_token\_minutes}) ;
  \item Contrôles de rôle sur DELETE et export CSV (vérification \texttt{role == admin}).
\end{itemize}

\section{Secrets et configuration}

\begin{table}[htbp]
  \centering
  \caption{Gestion des secrets}
  \label{tab:secrets}
  \begin{tabular}{@{}ll@{}}
    \toprule
    \textbf{Secret} & \textbf{Stockage} \\
    \midrule
    \texttt{SECRET\_JWT} & Secret K8s / \texttt{.env} (gitignoré) \\
    \texttt{DATABASE\_URL} & Secret K8s (mot de passe PG inclus) \\
    Mots de passe démo & Seed — à changer en prod \\
    Grafana admin & \texttt{TF\_VAR\_grafana\_admin\_password} \\
    \bottomrule
  \end{tabular}
\end{table}

\texttt{.env.example} documente les clés sans valeurs sensibles réelles. \texttt{.gitignore} exclut \texttt{.env}, \texttt{.kube/}, \texttt{*.tfstate}.

\section{CORS}

\texttt{CORS\_ORIGINS} liste les origines autorisées (frontend dev sur :3000, cluster sur \texttt{http://localhost}). En production, la liste serait restreinte au domaine client réel.

\section{Réseau Kubernetes}

Les Services clusterIP ne sont pas exposés sauf via Ingress (80/443). LocalStack et PostgreSQL ne sont pas publiés sur l'hôte en mode cluster (sauf port-forward manuel pour debug).

\section{Tests automatisés}

\subsection{Backend (pytest)}

Quatre tests couvrent :
\begin{itemize}
  \item connexion + liste produits ;
  \item création produit ;
  \item entrée de stock ;
  \item sortie refusée si stock insuffisant (\texttt{stock\_insuffisant}).
\end{itemize}

Fixtures : client FastAPI, jeton admin, base SQLite mémoire, override de \texttt{obtenir\_session}.

Le fichier \texttt{conftest.py} isole chaque test :
\begin{lstlisting}[language=Python,caption={Stratégie de test (extrait conftest)}]
os.environ["MODE_TEST"] = "1"
# SQLite :memory: + create_all avant chaque fonction de test
@pytest.fixture
def client():
    with TestClient(application) as c:
        yield c
\end{lstlisting}

Cette approche garantit un état vierge par test sans dépendre d'un Postgres déjà peuplé par le seed de production.

\subsection{Frontend (Vitest)}

Test de rendu de la page connexion avec providers Auth/Toast et MemoryRouter.

\subsection{Smoke tests (\texttt{make verify})}

\texttt{scripts/verifier.sh} enchaîne :
\begin{itemize}
  \item mode \texttt{dev} ou \texttt{cluster} ;
  \item curl sur \texttt{/health}, login, liste produits, page web ;
  \item en cluster : vérification pods Running.
\end{itemize}

\section{Analyse statique}

\begin{itemize}
  \item \textbf{ruff} — backend Python (CI + \texttt{make lint}) ;
  \item \textbf{eslint} — frontend TypeScript/React (CI).
\end{itemize}

\section{Pistes de durcissement production}

\begin{itemize}
  \item TLS terminé par cert-manager + Let's Encrypt ;
  \item rotation des secrets via External Secrets Operator ;
  \item politiques NetworkPolicy limitant le trafic inter-pods ;
  \item rate limiting sur l'Ingress ;
  \item audit des dépendances (\texttt{pip audit}, \texttt{npm audit}).
\end{itemize}
